%======================================================================
% 附录L：分章自测题库（含答案与提示）
%======================================================================
% 章标题 \appchapter{附录L：...} 由主文件发出，本文件只提供正文。
% 设计约定：每章固定 6 题 = 2 概念 + 2 推导 + 2 估算；每题给「答案要点」。
% 版权约定：只借鉴经典教材的题型与难度分级，题目、数值由本书自拟并独立复算。

\newcommand{\tkans}[1]{\par\vspace{0.5pt}{\small\textbf{答案：}#1}\par}
\newcommand{\tksrc}[1]{{\small\textbf{取材：}#1}\par\vspace{1pt}}

本附录给出一份\textbf{可执行的分章自测题库}：第 1--12 章每章 6 题，固定为 \textbf{2 道概念题 + 2 道推导题 + 2 道估算题}，每题附「答案要点」（给思路与关键数字，而不是只给一个数）。

它与另外两个附录的分工是：附录 K 是\textbf{范例}，给八道跨教材经典题与完整解答，用来学「怎么写」。附录 L 是\textbf{自测}，先合书独立做再看要点，用来查「会不会」。附录 I 是\textbf{策略}，管怎么安排复习、怎么处理错题。

\tksrc{每节的题型与难度分级参照 Chen 与 Fitzpatrick 教材对应章的习题类型（逐章对照见附录 J 的表 \ref{tab:app-j-map} 与本节末的来源表）；\textbf{所有题目与数值均由本书自拟，并已用脚本独立复算}，不复制任何原书题目。}

\begin{tip}{题库用法}{tip-appl-howto}
用法建议：先按附录 I 的「合书复述」独立做完一章的 6 题，再对答案要点。若一题的结果与文中\textbf{差一个量级}，先别怀疑公式，按附录 F 的顺序排查：温度用错（电子 / 离子、eV / K）、漏掉平方根、SI 与 CGS 混用、把 $kT$ 与 $T$ 等同。数值题一律给到两位有效数字即可，重点是量级与结论方向。
\end{tip}


\section{第 1 章：数学与统计力学预备知识}
\tksrc{Chen 第 1 章与附录 A（单位、矢量关系）、Fitzpatrick 各章散见的数学工具。}

\begin{enumerate}
\item （概念）德拜屏蔽与高斯定理矛盾吗？为什么屏蔽后的势不再是库仑势？
\tkans{不矛盾。高斯定理对\textbf{总场}成立：测试电荷周围的反号屏蔽云贡献的通量与点电荷通量部分抵消，$r\gg\lambda_D$ 时净通量趋于零。库仑势只是无屏蔽极限（$N_D\to\infty$）；德拜势 $(\frac{q}{4\pi\eps_0 r})e^{-r/\lambda_D}$ 是把屏蔽云的平均场计入后的结果。}

\item （概念）写出麦克斯韦速率分布的三个前提，并说明为什么最概然速率、平均速率、均方根速率三者不相等。
\tkans{前提：无外场、热平衡、非简并（经典统计）。速率分布 $f(v)\propto v^2e^{-mv^2/2k_BT}$ 不是关于最概然值的对称分布，$v^2$ 因子把分布向高速侧拉长，因此 $v_{th}:v_{avg}:v_{rms}=1:1.128:1.225$（三者依次增大）。}

\item （推导）由麦克斯韦分布导出 $v_{th}=\sqrt{2k_BT/m}$ 与 $v_{avg}=\sqrt{8k_BT/\pi m}$，并计算 $T_e=10\ \mathrm{eV}$ 时电子的两个数值。
\tkans{$v_{th}$ 由 $f(v)$ 的极大值给出，$v_{avg}$ 由分布的一阶矩给出，二者不相等。电子：$v_{th}=\sqrt{2T_e/m_e}=1.9\times10^{6}\ \mathrm{m/s}$，$v_{avg}=(2/\sqrt{\pi})\,v_{th}=2.1\times10^{6}\ \mathrm{m/s}$。注意两者相差 13\%，口算时不要混用。}

\item （推导）证明理想气体 $\langle\frac12mv^2\rangle=\frac32k_BT$，并说明「$1\ \mathrm{eV}$ 对应多少开尔文」。
\tkans{三个平动自由度各贡献 $\frac12k_BT$（能量均分）。由 $k_BT=e\cdot1\ \mathrm{V}$ 得 $1\ \mathrm{eV}=e/k_B=1.602\times10^{-19}/1.381\times10^{-23}=11605\ \mathrm{K}$。聚变等离子体「$10\ \mathrm{keV}$」即 $1.16\times10^{8}\ \mathrm{K}$，温度极高但粒子数密度低，压强仍在工程可及范围。}

\item （估算）$1\ \mathrm{Pa}$、$300\ \mathrm{K}$ 的中性气体数密度是多少？与 $n=10^{18}\ \mathrm{m^{-3}}$ 的「等离子体」相比，真空度谁更高？
\tkans{$n=p/k_BT=1/(1.381\times10^{-23}\times300)=2.4\times10^{20}\ \mathrm{m^{-3}}$，比 $10^{18}$ 高 240 倍。也就是说 $10^{18}\ \mathrm{m^{-3}}$ 的等离子体所处的本底真空必须优于 $10^{-2}\ \mathrm{Pa}$ 量级，否则中性粒子碰撞会主导，这正是实验室装置要抽高真空的原因。}

\item （估算）若 $n=10^{15}\ \mathrm{m^{-3}}$、$T=1\ \mathrm{eV}$，求 $\lambda_D$ 与德拜球内粒子数 $N_D$，并判断统计描述是否成立。
\tkans{$\lambda_D=7430\sqrt{T[\mathrm{eV}]/n[\mathrm{m^{-3}}]}=7430\sqrt{10^{-15}}=2.4\times10^{-4}\ \mathrm{m}$（0.24 mm）；$N_D=\frac{4\pi}{3}n\lambda_D^3\approx5.4\times10^{4}\gg1$。统计描述成立，但 $\lambda_D$ 已到毫米量级，装置尺度必须远大于 0.24 mm 才能谈「等离子体」（判据 $\lambda_D\ll L$ 与 $N_D\gg1$ 缺一不可）。}
\end{enumerate}


\section{第 2 章：什么是等离子体}
\tksrc{Chen 第 1 章（判据与参数）、Fitzpatrick 第 1 章。}

\begin{enumerate}
\item （概念）等离子体的三条判据各解决什么问题？
\tkans{$\lambda_D\ll L$：保证准中性、外场能被屏蔽（「集体」的前提）；$N_D\gg1$：保证德拜球内有足够多粒子，统计（连续介质）描述成立；$\omega_{pe}\tau>1$（$\tau$ 为碰撞时间）：保证集体振荡在两次碰撞之间能完成，等离子体行为不被碰撞抹平。三者分别对应空间尺度、统计性、时间尺度。}

\item （概念）为什么弱电离的辉光放电和完全电离的日冕都算等离子体，而火焰和电离层 D 层边缘要谨慎？
\tkans{判据看的是\textbf{集体性}而不是电离度：辉光放电 $n$ 大、$T_e$ 高，$\lambda_D$ 小、$N_D$ 大，仍满足三条判据；日冕虽稀薄，但尺度 $L$ 极大，$\lambda_D\ll L$ 成立。火焰电离度极低、且以化学反应与对流为主，集体行为被淹没；D 层边缘 $n$ 很低而碰撞频繁（$\omega_{pe}\tau<1$），也不满足判据。}

\item （推导）由泊松方程与玻尔兹曼因子 $n_e=n_0e^{e\phi/k_BT_e}$ 导出 $\lambda_D=\sqrt{\eps_0k_BT_e/(ne^2)}$。
\tkans{$\nabla^2\phi=-e(n_i-n_e)/\eps_0$；对 $|e\phi|\ll k_BT_e$ 展开 $n_e\approx n_0(1+e\phi/k_BT_e)$，代入得 $\nabla^2\phi=\phi/\lambda_D^2$，故 $\lambda_D=\sqrt{\eps_0k_BT_e/(n_0e^2)}$。工程式：$\lambda_D=7430\sqrt{T_e[\mathrm{eV}]/n[\mathrm{m^{-3}}]}\ \mathrm{m}$。}

\item （推导）证明德拜势在 $r=\lambda_D$ 处衰减到裸库仑势的 $1/e$，并求 $r=5\lambda_D$ 处剩余的百分比。
\tkans{$\phi\propto e^{-r/\lambda_D}/r$，故 $r=\lambda_D$ 时因子为 $e^{-1}=36.8\%$；$r=5\lambda_D$ 时为 $e^{-5}=0.67\%$，五个德拜长度之外屏蔽已经非常彻底，这就是「等离子体把电场关在自己家里」的定量含义。}

\item （估算）荧光灯正柱 $n=10^{16}\ \mathrm{m^{-3}}$、$kT_e=1\ \mathrm{eV}$：求 $\lambda_D$ 与 $N_D$，并判断它是否满足等离子体判据。
\tkans{$\lambda_D=7430\sqrt{1/10^{16}}=7.4\times10^{-5}\ \mathrm{m}$（74 $\mathrm{\mu m}$）；$N_D=\frac{4\pi}{3}\times10^{16}\times(7.43\times10^{-5})^3\approx1.7\times10^{4}\gg1$。管径量级为厘米，$\lambda_D\ll L$；电子温度 $1\ \mathrm{eV}$ 而气体接近室温，正是低温等离子体「电子热、离子冷」的典型参数。}

\item （估算）托卡马克边界（$n=10^{19}\ \mathrm{m^{-3}}$、$T_e=100\ \mathrm{eV}$）与日冕（$n=10^{15}\ \mathrm{m^{-3}}$、$T=100\ \mathrm{eV}$）的 $\lambda_D$ 之比是多少？
\tkans{$\lambda_D\propto\sqrt{T/n}$。边界：$\lambda_D=2.4\times10^{-5}\ \mathrm{m}$；日冕：$2.4\times10^{-3}\ \mathrm{m}$。比值 $\sqrt{10^{4}}=100$，即日冕的德拜长度比托卡马克边界大两个量级，密度低四个量级、温度相同。这解释了为什么日冕能容纳上百公里尺度的电流片。}
\end{enumerate}


\section{第 3 章：单粒子运动}
\tksrc{Chen 第 2 章、Fitzpatrick 第 2 章（平均化、绝热不变量）。}

\begin{enumerate}
\item （概念）$\vect{E}\times\vect{B}$ 漂移为什么与电荷符号和质量都无关？这会带来什么后果？
\tkans{运动方程 $m\diff\vect{v}/\diff t=q(\vect{E}+\vect{v}\times\vect{B})$ 的漂移解为 $\vect{v}_E=\vect{E}\times\vect{B}/B^2$，$q$ 与 $m$ 在求解过程中被消掉。后果：电子与离子\textbf{同向同速}漂移，不产生电荷分离，但整个等离子体一起平移，这是等离子体输运与「$E\times B$ 搅拌」的根源。}

\item （概念）磁镜为什么能反射粒子？损失锥的临界角由什么决定？
\tkans{磁矩 $\mu=mv_\perp^2/2B$ 是绝热不变量：$B$ 增大时 $v_\perp$ 增大，由能量守恒 $v_\parallel$ 减小，最终 $v_\parallel=0$ 而被反射。临界投掷角满足 $\sin^2\theta_c=B_{\min}/B_{\max}=1/R_m$（$R_m$ 为镜比），投掷角小于 $\theta_c$ 的粒子落入损失锥并逃逸。}

\item （推导）由 $m\diff\vect{v}/\diff t=q\vect{v}\times\vect{B}$（$\vect{B}=B\hat{z}$）导出 $\omega_c=|q|B/m$ 与 $\rho_L=v_\perp/\omega_c$，并说明回旋方向由什么决定。
\tkans{取 $\vect{v}=(v_\perp\cos\omega_ct,\,-v_\perp\sin\omega_ct,0)$ 代入即得圆周运动，角频率 $\omega_c=|q|B/m$，半径 $\rho_L=v_\perp/\omega_c=mv_\perp/|q|B$。旋转的\textbf{手性}（正电荷逆时针、负电荷顺时针，沿 $\vect{B}$ 方向看）由 $q$ 的符号决定，而频率与半径只依赖 $|q|/m$。}

\item （推导）托卡马克 $B=5\ \mathrm{T}$、氘离子 $T_i=10\ \mathrm{keV}$：求 $\omega_{ci}$、$v_i$ 与 $\rho_{Li}$。
\tkans{$\omega_{ci}=eB/m_i=1.602\times10^{-19}\times5/3.344\times10^{-27}=2.4\times10^{8}\ \mathrm{rad/s}$（$f_{ci}=38\ \mathrm{MHz}$）；$v_i=\sqrt{2k_BT_i/m_i}=9.8\times10^{5}\ \mathrm{m/s}$；$\rho_{Li}=v_i/\omega_{ci}=4.1\times10^{-3}\ \mathrm{m}$（4.1 mm）。离子回旋半径远小于装置半径（约 1--2 m），是磁约束可行的基本前提。}

\item （估算）地球磁层 $B=0.5\ \mathrm{G}$、$1\ \mathrm{keV}$ 质子：回旋半径是多少？为什么磁层能「困住」太阳风粒子？
\tkans{$\omega_{cp}=eB/m_p=1.602\times10^{-19}\times5\times10^{-5}/1.673\times10^{-27}=4.8\times10^{3}\ \mathrm{rad/s}$；$v=\sqrt{2k_BT/m_p}=4.4\times10^{5}\ \mathrm{m/s}$；$\rho_L=4.4\times10^{5}/4.8\times10^{3}=91\ \mathrm{m}$。它比磁层尺度（$10^{7}\ \mathrm{m}$）小五个量级，粒子只能沿磁力线螺旋前进，被磁场「串」在磁力线上，这就是辐射带与极光沉降的几何基础。}

\item （估算）$B=1\ \mathrm{T}$、$T_e=100\ \mathrm{eV}$ 的电子：回旋频率 $f_{ce}$ 与回旋半径 $\rho_{Le}$ 各是多少？与 $1\ \mathrm{cm}$ 探针尺度相比说明了什么？
\tkans{$\omega_{ce}=eB/m_e=1.76\times10^{11}\ \mathrm{rad/s}$，$f_{ce}=28\ \mathrm{GHz}$（毫米波段，可用于电子回旋加热与发射诊断）；$v_{te}=5.9\times10^{6}\ \mathrm{m/s}$，$\rho_{Le}=3.4\times10^{-5}\ \mathrm{m}$（34 $\mathrm{\mu m}$）。$\rho_{Le}\ll1\ \mathrm{cm}$：电子被强烈磁化，其横越磁力线的运动被冻结，只有沿磁力线才自由，这是磁化等离子体各向异性输运的起点。}
\end{enumerate}

\section{第 4 章：流体描述基础}
\tksrc{Chen 第 3 章、Fitzpatrick 第 4 章（Chapman--Enskog、Braginskii 输运）。}

\begin{enumerate}
\item （概念）连续性方程与动量方程各表达什么守恒？双流体模型还需要什么才能闭合？
\tkans{连续性方程 $\partial_tn+\nabla\cdot(n\vect{u})=0$ 是粒子数（质量）守恒；动量方程是动量守恒，含压强梯度、洛伦兹力 $\vect{j}\times\vect{B}$ 与碰撞摩擦。未知量多于方程，必须补\textbf{状态方程或能量方程}（等温 / 绝热假设）并加上麦克斯韦方程组，才能闭合。}

\item （概念）抗磁漂移的物理来源是什么？为什么它与电荷符号有关，而 $\vect{E}\times\vect{B}$ 漂移无关？
\tkans{压强梯度使回旋轨道两侧的受力不平衡，产生 $\vect{v}_D=-(\nabla p\times\vect{B})/(qnB^2)$：式中含 $q$，因此电子与离子\textbf{反向}漂移，形成抗磁电流（这正是等离子体电流与磁场的「自带」来源）。$\vect{E}\times\vect{B}$ 漂移不含 $q$，故与电荷符号无关。}

\item （推导）由流体动量方程导出抗磁漂移的标量形式 $v_{dia}=k_BT/(qBL_n)$（$L_n=|\nabla n/n|^{-1}$）。
\tkans{对稳态、无碰撞、各向同性压强，动量方程 $0=-\nabla p+qn(\vect{v}\times\vect{B})$ 给出 $\vect{v}_\perp=(\nabla p\times\vect{B})/(qnB^2)$。取 $\nabla p=k_BT\nabla n$、$|\nabla n|/n=1/L_n$，得 $v_{dia}=k_BT/(qBL_n)$；方向垂直于 $\nabla n$ 与 $\vect{B}$，电子与离子相反。}

\item （推导）$n=10^{19}\ \mathrm{m^{-3}}$、$T=100\ \mathrm{eV}$、$B=2\ \mathrm{T}$、$L_n=0.1\ \mathrm{m}$：求 $v_{dia}$，以及由此产生的抗磁电流密度 $j$。
\tkans{$v_{dia}=k_BT/(eBL_n)=1.602\times10^{-17}/(1.602\times10^{-19}\times2\times0.1)=500\ \mathrm{m/s}$。离子与电子速率相同、方向相反，故 $j=ne(v_i-v_e)=10^{19}\times1.602\times10^{-19}\times1000=1.6\times10^{3}\ \mathrm{A/m^2}$。1 cm 的密度梯度就能驱动千安每平方米量级的电流。}

\item （估算）氘等离子体 $T_e=10\ \mathrm{eV}$ 的离子声速 $c_s=\sqrt{k_BT_e/m_i}$ 是多少？与 $1\ \mathrm{eV}$ 的氩等离子体相比呢？
\tkans{氘（$m_i=3.34\times10^{-27}\ \mathrm{kg}$）：$c_s=\sqrt{T_e/m_i}=2.2\times10^{4}\ \mathrm{m/s}$；氩（$m_i=39.95\,m_p$）：$c_s=1.6\times10^{3}\ \mathrm{m/s}$。相差 14 倍，完全来自 $1/\sqrt{m_i}$，重离子等离子体的特征速度慢得多（鞘层、玻姆判据与预鞘加速都按这个速度标定）。}

\item （估算）若碰撞平均自由程 $\lambda_{mfp}=10^{-2}\ \mathrm{m}$、装置尺度 $L=0.5\ \mathrm{m}$，流体描述是否成立？
\tkans{$\lambda_{mfp}/L=0.02\ll1$：碰撞足够频繁，速度分布接近局域麦克斯韦分布，流体（局域热力学平衡）近似可用。反之若 $\lambda_{mfp}\gtrsim L$（如聚变芯部、空间等离子体），必须用动理学描述，这正是第 7 章与第 8 章的分工。}
\end{enumerate}


\section{第 5 章：磁流体力学}
\tksrc{Chen 第 3、6 章、Fitzpatrick 第 8 章（磁压、磁冻结、MHD 波），进阶参照 Freidberg \textit{Ideal MHD}。}

\begin{enumerate}
\item （概念）理想 MHD 的三条基本假设与两个「冻结」结论是什么？
\tkans{假设：准中性（$L\gg\lambda_D$）、单流体（离子与电子一起运动，$L\gg\rho_i$、$\omega\ll\omega_{ci}$）、理想欧姆定律 $\vect{E}+\vect{v}\times\vect{B}=0$（等价于 $R_m\gg1$）。结论：磁通量冻结（穿过任一流体回路的磁通守恒）、磁场与流体「焊」在一起；等价说法是磁力线像被流体拖着的橡皮筋。}

\item （概念）为什么说 $\vect{j}\times\vect{B}$ 可以写成「磁张力 + 磁压」？
\tkans{用矢量恒等式 $\vect{j}\times\vect{B}=(\vect{B}\cdot\nabla)\vect{B}/\mu_0-\nabla(B^2/2\mu_0)$。第一项沿磁力线方向，是\textbf{磁张力}（把弯曲的磁力线拉直）；第二项是各向同性的\textbf{磁压}（$B^2/2\mu_0$，把磁力线彼此推开）。于是 MHD 动量方程获得直观图像：磁力线是绷紧的弦，磁场是能推的气体。}

\item （推导）由感应方程 $\partial_t\vect{B}=\nabla\times(\vect{v}\times\vect{B})+(\eta/\mu_0)\nabla^2\vect{B}$ 导出磁扩散时间 $\tau_\eta=\mu_0L^2/\eta$ 与磁雷诺数 $R_m=\mu_0Lv/\eta$。
\tkans{忽略对流项，方程化为扩散方程，特征时间 $\tau_\eta=\mu_0L^2/\eta$；两项之比即 $R_m=\tau_\eta/\tau_{flow}=\mu_0Lv/\eta$。$R_m\gg1$ 时对流项主导，冻结成立；$R_m\ll1$ 时磁场自由扩散。}

\item （推导）$n=10^{20}\ \mathrm{m^{-3}}$、$B=5\ \mathrm{T}$ 的氘等离子体：求 Alfvén 速度，并用 $v_A/c$ 说明非相对论性。
\tkans{$v_A=B/\sqrt{\mu_0nm_i}=5/\sqrt{1.2566\times10^{-6}\times10^{20}\times3.344\times10^{-27}}=5/6.48\times10^{-7}=7.7\times10^{6}\ \mathrm{m/s}$；$v_A/c=2.6\%$。$v_A/c\ll1$ 只说明\textbf{非相对论 Alfv\'en 公式适用}（位移电流可以忽略），\textbf{不能据此断定磁冻结成立}：磁冻结要求理想欧姆定律成立，至少要 $R_m=\mu_0Lv/\eta\gg1$，还要检查霍尔项与电子压强项（正文第 5 章「MHD 适用条件：三道独立的门槛，别只看一道」）。本题未给温度、电阻率与流动尺度，无法判断（对聚变参数，$v_A$ 通常为 $10^{6}$--$10^{7}\ \mathrm{m/s}$）。}

\item （估算）$T_e=10\ \mathrm{eV}$、$\ln\Lambda=10$ 的氘等离子体：$\eta_\parallel\approx5.2\times10^{-5}\ln\Lambda\,T_e[\mathrm{eV}]^{-3/2}$，求 $\eta_\parallel$、$L=1\ \mathrm{m}$ 的 $\tau_\eta$，以及 $v=10^{4}\ \mathrm{m/s}$ 时的 $R_m$。磁冻结在这里成立吗？
\tkans{$\eta=5.2\times10^{-5}\times10/10^{1.5}=1.6\times10^{-5}\ \Omega\cdot\mathrm{m}$；$\tau_\eta=\mu_0L^2/\eta=1.2566\times10^{-6}/1.64\times10^{-5}=0.076\ \mathrm{s}$；$R_m=\mu_0Lv/\eta=1.2566\times10^{-6}\times1\times10^{4}/1.64\times10^{-5}=7.6\times10^{2}$。$R_m\gg1$ 成立，但只有两个多量级，若把 $L$ 缩到 1 cm 或把速度降到 $10^{3}\ \mathrm{m/s}$，$R_m$ 就跌到 10 以下，「理想」随即失效。结论：磁冻结是\textbf{需要检验}的近似，不是自动成立的定理。}

\item （估算）$\beta=2\mu_0nk_BT/B^2$：$n=10^{20}\ \mathrm{m^{-3}}$、$T=10\ \mathrm{keV}$、$B=5.3\ \mathrm{T}$ 时 $\beta$ 是多少？若要把 $\beta$ 压到 3\%（托卡马克运行区），需要多大磁场？
\tkans{$p=nk_BT=10^{20}\times10^{4}\times1.602\times10^{-19}=1.6\times10^{5}\ \mathrm{Pa}$；$\beta=2\times1.2566\times10^{-6}\times1.602\times10^{5}/5.3^2=1.4\%$。要 $\beta=3\%$：$B=\sqrt{2\mu_0p/0.03}=\sqrt{0.4026/0.03}=3.7\ \mathrm{T}$。可见同温度下提高磁场是压低 $\beta$（换取 MHD 稳定裕度）的直接手段。}
\end{enumerate}


\section{第 6 章：等离子体中的波}
\tksrc{Chen 第 4 章、Fitzpatrick 第 5、6、7 章（冷、非均匀、温等离子体），进阶参照 Stix \textit{Waves in Plasmas}。}

\begin{enumerate}
\item （概念）Bohm--Gross 色散 $\omega^2=\omega_{pe}^2+3k^2v_{te}^2$ 中两项各代表什么？$k\lambda_D\ll1$ 时回到什么？
\tkans{$\omega_{pe}^2$ 是电荷分离产生的静电恢复力（冷等离子体振荡）；$3k^2v_{te}^2$ 是电子热压强提供的恢复力（把 $3v_{te}^2$ 换成 $\gamma k_BT_e/m_e$ 即一维绝热声波）。$k\lambda_D\ll1$ 时热项可忽略，$\omega\to\omega_{pe}$：长波极限下电子集体振荡频率与 $k$ 无关。}

\item （概念）截止与共振有什么区别？为什么反射计能测量密度剖面？
\tkans{截止：$n$ 降到使 $k^2=0$（$\omega=\omega_{pe}$），波不能传播而被反射；共振：$k\to\infty$（如 $\omega=\omega_{ce}$ 附近的共振层），波被强烈吸收。反射计向等离子体发射扫频波，测到反射时的频率即给出截止层密度 $n_c(\omega)$；由不同频率的反射时间（或相位）反演出密度剖面。}

\item （推导）由 $\omega^2=\omega_{pe}^2+k^2c^2$ 导出截止密度 $n_c=\eps_0m_e\omega^2/e^2$，并求 $f=2.45\ \mathrm{GHz}$ 时的 $n_c$。
\tkans{令 $k=0$ 得 $\omega^2=\omega_{pe}^2=n_ce^2/(\eps_0m_e)$，故 $n_c=\eps_0m_e\omega^2/e^2$。$f=2.45\ \mathrm{GHz}$：$\omega=1.539\times10^{10}\ \mathrm{rad/s}$，$n_c=8.854\times10^{-12}\times9.109\times10^{-31}\times2.369\times10^{20}/2.566\times10^{-38}=7.4\times10^{16}\ \mathrm{m^{-3}}$。微波炉频率对应的正是这个量级的密度，等离子体工艺中「过密」与「欠密」的分界。}

\item （推导）离子声波 $\omega=kc_s$、$c_s=\sqrt{k_BT_e/m_i}$ 在什么条件下成立？求 $T_e=10\ \mathrm{eV}$ 时氘等离子体的 $c_s$。
\tkans{条件：$T_e\gg T_i$（离子被电子压强推着走）、$k\lambda_D\ll1$（准中性）。$c_s=\sqrt{T_e/m_i}=2.2\times10^{4}\ \mathrm{m/s}$，与第 4 章估算题一致；这个速度同时是玻姆速度与预鞘加速的标定速度。}

\item （估算）$n=10^{18}\ \mathrm{m^{-3}}$ 的等离子体频率 $f_{pe}$ 是多少？$2.45\ \mathrm{GHz}$ 的微波能穿透吗？
\tkans{$\omega_{pe}=\sqrt{ne^2/(\eps_0m_e)}=5.6\times10^{10}\ \mathrm{rad/s}$，$f_{pe}=9.0\ \mathrm{GHz}$。$f_{pe}>2.45\ \mathrm{GHz}$，即 $\omega<\omega_{pe}$，波是渐逝的，微波被反射（该等离子体对微波「过密」）。这正是「微波不能穿透高密度等离子体」的定量判据，也是 2.45 GHz 源难以维持高密度放电的原因。}

\item （估算）$n=10^{20}\ \mathrm{m^{-3}}$、$B=5\ \mathrm{T}$ 的氘等离子体中，Alfvén 波沿 $1\ \mathrm{m}$ 传播需要多久？
\tkans{$v_A=7.7\times10^{6}\ \mathrm{m/s}$（第 5 章已算），$t=L/v_A=1.3\times10^{-7}\ \mathrm{s}=0.13\ \mathrm{\mu s}$。微秒量级，比 MHD 不稳定性的增长时间快得多，说明理想 MHD 波可以在扰动增长期间多次往返，磁力线的「弹性」响应是瞬时的。}
\end{enumerate}

\section{第 7 章：动理学理论}
\tksrc{Chen 第 7 章、Fitzpatrick 第 3、7 章（Landau 阻尼、速度空间不稳定性），进阶参照 Krall \& Trivelpiece。}

\begin{enumerate}
\item （概念）弗拉索夫方程与玻尔兹曼方程差在哪里？为什么「无碰撞」也能有阻尼？
\tkans{弗拉索夫方程 $\partial_tf+\vect{v}\cdot\nabla f+(q/m)(\vect{E}+\vect{v}\times\vect{B})\cdot\nabla_vf=0$ 只保留自洽平均场，丢掉碰撞积分；玻尔兹曼方程多一项 $(\delta f/\delta t)_c$。Landau 阻尼不是碰撞阻尼，而是\textbf{波与共振粒子的可逆能量交换}：相速度附近的粒子被波加速或被减速，净效果是把相干波能转移为粒子的无规动能，因此熵增、波衰减，而无需任何碰撞。}

\item （概念）共振条件 $\omega-kv_\parallel=0$ 意味着什么？什么样的粒子主导 Landau 阻尼？
\tkans{平行速度等于波相速度 $v_\phi=\omega/k$ 的粒子与波「同相」共振，能持续交换能量。主导者是分布函数在 $v_\phi$ 处的斜率 $\partial f_0/\partial v$：斜率越平缓（或越靠尾部）贡献越小。$v_\phi\sim v_{te}$ 时阻尼最强；$v_\phi$ 远在分布尾部之外时阻尼指数小，这解释了为什么快波（$v_\phi\gg v_{te}$）几乎无阻尼。}

\item （推导）从弗拉索夫方程线性化（$f=f_0+\delta f$，$\delta f\propto e^{i(kx-\omega t)}$）导出 $\delta f=\dfrac{qE}{m}\dfrac{\partial f_0/\partial v}{i(kv-\omega)}$，并说明分母的极点如何处理。
\tkans{代入并保留一阶项得 $(-i\omega+ikv)\delta f+(qE/m)\partial_vf_0=0$，即上式。极点 $v=\omega/k$ 落在积分围道上，物理上对应共振粒子：绕开极点的方式（朗道围道，$\omega\to\omega+i0^+$）决定了积分取主值加半留数，其虚部给出阻尼/增长，极点的处理方式就是阻尼率的来源。}

\item （推导）写出麦克斯韦分布下的弱阻尼公式，并计算 $k\lambda_D=0.3$ 时的 $\gamma/\omega$。
\tkans{$\dfrac{\gamma}{\omega}=-\sqrt{\dfrac{\pi}{8}}\dfrac{1}{(k\lambda_D)^3}\exp\!\left[-\dfrac{1}{2k^2\lambda_D^2}-\dfrac{3}{2}\right]$。$k\lambda_D=0.3$：$(k\lambda_D)^{-3}=37.04$，指数 $\exp[-5.556-1.5]=e^{-7.056}=8.64\times10^{-4}$，故 $\gamma/\omega=-0.6266\times37.04\times8.64\times10^{-4}=-2.0\times10^{-2}$，即每周期衰减约 12\%。}

\item （估算）$k\lambda_D=0.25$ 与 $k\lambda_D=0.2$ 的 $\gamma/\omega$ 各是多少？为什么两者差近两个量级？
\tkans{$k\lambda_D=0.25$：$\gamma/\omega=-3.0\times10^{-3}$（每周期约 1.9\%）；$k\lambda_D=0.2$：$-6.5\times10^{-5}$（0.04\%）。差别几乎全部来自指数因子 $e^{-1/2k^2\lambda_D^2}$：从 $e^{-8}$ 降到 $e^{-12.5}$，倍率 $e^{-4.5}=0.011$。这就是「长波几乎无阻尼、$k\lambda_D\sim0.3$ 才明显」的定量含义，也是为什么必须把波长与 $\lambda_D$ 一起报告。}

\item （估算）双流不稳定性：$n=10^{18}\ \mathrm{m^{-3}}$、$\gamma_{\max}=\omega_{pe}/\sqrt{8}$，求 $\omega_{pe}$、$\gamma_{\max}$ 与特征增长时间。
\tkans{$\omega_{pe}=5.6\times10^{10}\ \mathrm{rad/s}$；$\gamma_{\max}=5.64\times10^{10}/2.828=2.0\times10^{10}\ \mathrm{s^{-1}}$；增长时间 $1/\gamma_{\max}=5.0\times10^{-11}\ \mathrm{s}=50\ \mathrm{ps}$。比任何碰撞或输运时间都快几个量级，这解释了为什么束-等离子体系统一旦满足条件就迅速进入强湍流（第 8 章「反常电阻」的动理学根源）。}
\end{enumerate}


\section{第 8 章：碰撞与输运}
\tksrc{Chen 第 5 章、Fitzpatrick 第 3、4 章（库仑对数、Rosenbluth 势），进阶参照 Helander \& Sigmar（新经典输运）。}

\begin{enumerate}
\item （概念）库仑对数 $\ln\Lambda$ 的物理意义是什么？为什么它几乎出现在所有碰撞率公式里？
\tkans{$\ln\Lambda=\ln(\lambda_D/b_{90})$，其中 $b_{90}$ 是使偏转角达 $90^\circ$ 的碰撞参数。库仑力长程，$90^\circ$ 散射的微分截面在小角发散，必须用\textbf{大量小角散射的累积}定义有效碰撞时间；积分给出对碰撞参数的对数权重 $\ln\Lambda$（典型 10--20，对 $n$、$T$ 只有弱依赖）。因此所有库仑碰撞率都带一个 $\ln\Lambda$ 因子，而它的不确定性（因子 2 以内）通常不影响结论。}

\item （概念）经典扩散与玻姆扩散的标度差异是什么？为什么实验上常观察到玻姆扩散？
\tkans{经典（碰撞）扩散 $D_\perp\propto n/T^{1/2}B^2$（对 $B$ 是 $1/B^2$），来自回旋运动被碰撞打断的随机行走；玻姆扩散 $D_B\propto k_BT_e/(eB)$（$1/B$），来自湍流涡的尺度 $\sim\rho_s$ 与寿命 $\sim1/\omega_*$。差异在 $B$ 的幂次：装置加大磁场时经典扩散迅速变小，而玻姆扩散只线性下降，因此在强磁场约束装置中湍流（玻姆类）输运常常主导。}

\item （推导）由 Spitzer 电阻率 $\eta_\parallel=5.2\times10^{-5}Z\ln\Lambda\,T_e[\mathrm{eV}]^{-3/2}\ \Omega\cdot\mathrm{m}$ 说明 $\eta\propto T^{-3/2}$，并计算温度从 $10\ \mathrm{eV}$ 升到 $1\ \mathrm{keV}$ 时 $\eta$ 下降多少倍。
\tkans{$\eta\propto T^{-3/2}$ 直接来自 $\eta\propto\nu_{ei}\propto nT^{-3/2}$（除以 $n$ 后仍留 $T^{-3/2}$）。倍率 $(1000/10)^{3/2}=100^{1.5}=1000$。温度升高 100 倍，电阻率降低 1000 倍，这是欧姆加热在高温下「自我熄火」的根源，也是为什么聚变装置必须靠中性束、射频等辅助加热。}

\item （推导）$n=10^{19}\ \mathrm{m^{-3}}$、$T_e=1\ \mathrm{keV}$、$\ln\Lambda=10$：求 $\eta_\parallel$ 与电子-离子碰撞时间 $\tau_{ei}$（$\nu_{ei}\approx2.9\times10^{-6}n\ln\Lambda\,T_e[\mathrm{eV}]^{-3/2}\ \mathrm{s^{-1}}$）。
\tkans{$\eta=5.2\times10^{-5}\times10/1000^{1.5}=1.6\times10^{-8}\ \Omega\cdot\mathrm{m}$；$\nu_{ei}=2.9\times10^{-6}\times10^{19}\times10/3.162\times10^{4}=9.2\times10^{9}\ \mathrm{s^{-1}}$，$\tau_{ei}=1/\nu_{ei}=1.1\times10^{-10}\ \mathrm{s}$。比较：$\omega_{pe}=5.6\times10^{9}\ \mathrm{rad/s}$，$\omega_{pe}\tau_{ei}=0.6$，量级为 1，说明该参数下集体效应与碰撞效应同等重要，「无碰撞」与「碰撞主导」都不是好的近似。}

\item （估算）上题电子的热速度 $v_{te}=1.9\times10^{7}\ \mathrm{m/s}$，平均自由程 $\lambda_{mfp}=v_{te}\tau_{ei}$ 是多少？与 $1\ \mathrm{m}$ 的装置尺度相比说明了什么？
\tkans{$\lambda_{mfp}=1.9\times10^{7}\times1.1\times10^{-10}=2.1\times10^{-3}\ \mathrm{m}$（2 mm），$\lambda_{mfp}/L=2\times10^{-3}\ll1$。虽然是「高温稀薄」等离子体，芯部的碰撞仍然足够频繁，局域输运系数（Spitzer、Braginskii）可用；但要注意这是\textbf{沿磁力线}的自由程，垂直方向的步长被回旋半径限制，两者相差 $\sim10^{4}$ 倍。}

\item （估算）中性粒子的电离/电荷交换截面取 $\sigma\approx10^{-19}\ \mathrm{m^2}$：边界 $n=10^{20}\ \mathrm{m^{-3}}$ 时穿透深度 $1/(n\sigma)$ 是多少？若边界密度升到 $10^{21}\ \mathrm{m^{-3}}$ 呢？
\tkans{$n=10^{20}$：$\lambda=1/(10^{20}\times10^{-19})=0.1\ \mathrm{m}$（10 cm）；$n=10^{21}$：$\lambda=1\ \mathrm{cm}$。穿透深度 $\propto1/n$，所以中性粒子（加料、再循环）只作用于边界薄层，芯部几乎完全电离，这也是偏滤器要把中性粒子「关」在靶板附近、以及中性束注入必须高能（才能边电离边穿入）的原因。}
\end{enumerate}


\section{第 9 章：等离子体不稳定性}
\tksrc{Chen 第 6、8 章、Fitzpatrick 第 7、9 章（撕裂模、重联），进阶参照 Freidberg 与 Wesson \textit{Tokamaks}。}

\begin{enumerate}
\item （概念）不稳定性的自由能来自哪里？「$\gamma>0$」的判据从哪里得到？
\tkans{自由能来源：密度/温度梯度（漂移波、交换模）、电流梯度（撕裂模、扭结模）、速度空间各向异性或束流（双流、损失锥不稳定性）、坏曲率与重力（交换模、瑞利--泰勒）。判据有两条路径：解色散关系看 $\omega$ 是否有正虚部（线性增长率），或用能量原理 $\delta W<0$（MHD 变分判据，不需要显式色散关系）。}

\item （概念）为什么「坏曲率」会驱动交换模？它与瑞利--泰勒不稳定性有什么类比？
\tkans{在曲率中心位于等离子体一侧（坏曲率）的区域，磁力线凹面朝向等离子体，磁压随半径下降得比等离子体压强快，等价于「重流体压在轻流体上」：把两根相邻磁通管交换位置会降低总能量，因此不稳定。类比关系是 $g_{eff}=c_s^2/R_c$ 扮演重力加速度，判据由 $\omega^2=-Akg_{eff}<0$ 给出，与流体瑞利--泰勒同形，差别只在「重力」来自磁力线弯曲而非真实引力。}

\item （推导）由 $\omega^2=-Akg$ 导出瑞利--泰勒增长率 $\gamma=\sqrt{Akg}$；再用 $g_{eff}=c_s^2/R_c$ 写出坏曲率交换模的增长率并说明稳定条件。
\tkans{$\omega^2<0$ 时 $\omega=\pm i\sqrt{Akg}$，故 $\gamma=\sqrt{Akg}$（$A$ 为 Atwood 数，$k$ 为波数）。交换模：$\gamma=\sqrt{kc_s^2/R_c}$。稳定条件：$k$ 与 $R_c$ 的组合使 $\gamma$ 被磁剪切或有限 $\beta$ 抑制，工程上表现为「坏曲率区足够小、磁剪切足够强」（把 $g_{eff}$ 变成沿磁力线的振荡平均，见磁井与平均最小 $B$）。}

\item （推导）漂移波色散 $\omega=\omega_*/(1+k_\perp^2\rho_s^2)$（$\omega_*=k_\perp\rho_sc_s/L_n$）：说明长波极限、增长率最大处的 $k_\perp$，以及湍流涡的尺度。
\tkans{$k_\perp\rho_s\ll1$ 时 $\omega\approx\omega_*$，频率正比于 $k_\perp$（沿抗磁方向传播）；$k_\perp\rho_s\gtrsim1$ 时离子有限回旋半径把涨落平均掉，$\omega$ 随 $k_\perp$ 下降，增长最强的尺度为 $k_\perp\rho_s\sim1$，即涡尺度 $\sim\rho_s$（离子声回旋半径，毫米量级）。这把「湍流输运」从定性图像落到一个可计算的尺度上。}

\item （估算）瑞利--泰勒：$A=1$、扰动波长 $6.3\ \mathrm{cm}$（$k=100\ \mathrm{m^{-1}}$）、$g=9.8\ \mathrm{m/s^2}$，求 $\gamma$ 与 e 折时间。
\tkans{$\gamma=\sqrt{Akg}=\sqrt{100\times9.8}=31\ \mathrm{s^{-1}}$，e 折时间 $1/\gamma=32\ \mathrm{ms}$。注意 $\gamma\propto\sqrt{k}$：波长越短的扰动长得越快（但受黏性/有限拉莫尔半径截断），这就是「界面扰动很快破碎成小尺度结构」的原因。}

\item （估算）托卡马克坏曲率区：$c_s=10^{5}\ \mathrm{m/s}$、$R_c=1\ \mathrm{m}$、$k=1\ \mathrm{m^{-1}}$，用 $\gamma=\sqrt{kc_s^2/R_c}$ 估算交换模增长率与增长时间。
\tkans{$g_{eff}=c_s^2/R_c=10^{10}\ \mathrm{m/s^2}$，$\gamma=\sqrt{kg_{eff}}=\sqrt{1.0\times10^{10}}=1.0\times10^{5}\ \mathrm{s^{-1}}$，增长时间 $10\ \mathrm{\mu s}$。比能量约束时间（$10^{-1}$--$10^{0}\ \mathrm{s}$）快四个量级：一旦坏曲率区与共振面重合，扰动会在远小于约束时间内长到非线性，所以托卡马克必须靠磁剪切与等离子体截面整形把「平均好曲率」做出来。}
\end{enumerate}

\section{第 10 章：约束与受控核聚变}
\tksrc{Chen 第 9、10 章、Fitzpatrick 第 8 章（太阳风、激波）；装置与工程量级参照 Goldston \& Rutherford 与 Wesson \textit{Tokamaks}。}

\begin{enumerate}
\item （概念）劳森判据的三个要素是什么？为什么写成「三乘积」而不是单一指标？
\tkans{密度 $n$、温度 $T$、能量约束时间 $\tau_E$。聚变功率密度 $\propto n^2\langle\sigma v\rangle$（两体过程），而单位体积热能 $\propto nT$、损失功率 $\propto nT/\tau_E$；两者之比要求 $nT\tau_E$ 足够大。三者的乘积形式说明\textbf{单靠提高密度或温度都不够}：还要把能量留住（$\tau_E$），这正是约束（磁场、惯性）而不是加热成为核心难点的原因。}

\item （概念）托卡马克与仿星器实现「旋转变换」的方式有何不同？各自的代价是什么？
\tkans{托卡马克靠\textbf{等离子体自身电流}产生极向场，从而让磁力线绕环螺旋（旋转变换）；仿星器完全用\textbf{外部三维线圈}的几何形状产生旋转变换。代价对比：托卡马克装置简单但需要驱动电流（射频/中性束），且电流剖面会驱动撕裂模、破裂等不稳定性；仿星器无净电流、无破裂风险，但线圈工程复杂、磁场误差敏感、输运优化困难。}

\item （推导）由零维功率平衡 $P_\alpha+P_{ext}=P_{loss}$ 与 $P_{fus}\propto n^2\langle\sigma v\rangle$ 说明点火条件如何化为 $nT\tau_E$ 的门槛，并给出 $Q$ 的定义。
\tkans{取 $P_\alpha=P_{fus}/5$（$\alpha$ 粒子带走 3.5 MeV/17.6 MeV），$P_{loss}=3nk_BT/\tau_E$。令 $P_{ext}=0$（点火）即得 $n^2\langle\sigma v\rangle E_\alpha/5=3nk_BT/\tau_E$，整理成 $nT\tau_E\ge\dfrac{15k_BT^2}{\langle\sigma v\rangle E_\alpha}$，右边只依赖 $T$，故门槛是三乘积。$Q=P_{fus}/P_{ext}$：$Q=1$ 得失相当，$Q=5$（ITER 目标）时 $\alpha$ 加热已占主导，$Q\to\infty$ 即点火。}

\item （推导）D--T 等离子体 $n_D=n_T=10^{20}\ \mathrm{m^{-3}}$、$\langle\sigma v\rangle=1.1\times10^{-22}\ \mathrm{m^3/s}$，每次反应释放 $17.6\ \mathrm{MeV}$：求聚变功率密度。
\tkans{$P_{fus}=n_Dn_T\langle\sigma v\rangle E_{fus}=10^{40}\times1.1\times10^{-22}\times(17.6\times10^{6}\times1.602\times10^{-19})=1.1\times10^{18}\times2.82\times10^{-12}=3.1\times10^{6}\ \mathrm{W/m^3}$，即约 $3\ \mathrm{MW/m^3}$。这是「堆芯量级」的功率密度，与裂变堆芯相当，这也是为什么聚变堆必须解决第一壁热流问题。}

\item （估算）接上题：$T=10\ \mathrm{keV}$ 时热能密度 $3nk_BT$ 是多少？$\alpha$ 加热功率密度 $P_\alpha$ 与点火门槛 $\tau_E$ 各是多少？
\tkans{$3nk_BT=3\times10^{20}\times10^{4}\times1.602\times10^{-19}=4.8\times10^{5}\ \mathrm{J/m^3}$；$P_\alpha=P_{fus}/5=6.2\times10^{5}\ \mathrm{W/m^3}$；自持要求 $\tau_E>3nk_BT/P_\alpha=0.77\ \mathrm{s}$。这是零维、忽略轫致辐射与杂质辐射的乐观估算；真实装置还要加上辐射损失与加热效率，门槛会更高（ITER 设计 $\tau_E\sim3\ \mathrm{s}$ 量级）。}

\item （估算）一个 $V=100\ \mathrm{m^3}$ 的 D--T 堆芯（参数同上）总聚变功率是多少？要达到 $1\ \mathrm{GW}$ 需要多大体积？
\tkans{$P=P_{fus}V=3.1\times10^{6}\times100=3.1\times10^{8}\ \mathrm{W}=310\ \mathrm{MW}$；要到 $1\ \mathrm{GW}$ 需 $V\approx320\ \mathrm{m^3}$（或把 $n$ 提高约 1.8 倍）。把「功率密度 × 体积」与工程尺度对上，是理解为什么聚变堆体积以千立方米计的最短路径。}
\end{enumerate}


\section{第 11 章：低温工业等离子体}
\tksrc{Chen 第 10 章、Fitzpatrick 第 4 章末（Langmuir 鞘层与探针），进阶参照 Lieberman \& Lichtenberg。}

\begin{enumerate}
\item （概念）低温等离子体为什么「电子热、离子冷」？这种非平衡状态带来了什么好处？
\tkans{电子质量小，在电场中被加速后与离子/中性粒子的弹性碰撞能量传递率正比于 $m_e/m_i$（$\sim10^{-4}$），因此电子从电场获得的能量难以传给重粒子：电子温度可达数个 eV，而气体温度接近室温。好处是\textbf{用低温衬底获得高能电子}，由它驱动电离、离解、激发与自由基生成，从而在材料表面实现刻蚀、沉积、改性（半导体工艺的基础）。}

\item （概念）鞘层为什么必然存在？浮动电位为什么是负的？
\tkans{等离子体与器壁接触时，电子热速度远大于离子玻姆速度（$v_{te}/v_B\sim\sqrt{m_i/m_e}\sim10^{2}$），初始电子通量远大于离子通量，器壁迅速充负电；负电位反过来排斥电子、吸引离子，直到两者通量相等。平衡电位即浮动电位，必然为负，量级为几个 $k_BT_e/e$。}

\item （推导）写出玻姆判据与 Child--Langmuir 定律，并给出由已知离子流密度求鞘层厚度的反解式。
\tkans{玻姆判据：离子进入鞘边时速度必须达到 $v_B=\sqrt{k_BT_e/m_i}$（否则鞘内离子密度下降快于电子，无法维持单调电位）。Child--Langmuir（无碰撞、空间电荷限制）$J_i=\dfrac{4\eps_0}{9}\sqrt{\dfrac{2e}{m_i}}\dfrac{|V_w|^{3/2}}{s^2}$，反解得 $s=\sqrt{\dfrac{4\eps_0}{9J_i}\sqrt{\dfrac{2e}{m_i}}|V_w|^{3/2}}$。鞘厚是 $\lambda_D$ 的若干倍，随 $|V_w|/T_e$ 的 $3/2$ 次方增长。}

\item （推导）氩等离子体 $T_e=3\ \mathrm{eV}$：求浮动电位 $V_f=\dfrac{k_BT_e}{2e}\ln\dfrac{m_i}{2\pi m_e}$ 与玻姆速度 $v_B$，并核对两者量级（$\sim5T_e/e$）。
\tkans{$\ln\dfrac{6.634\times10^{-26}}{2\pi\times9.109\times10^{-31}}=\ln(1.16\times10^{4})=9.36$，$V_f=1.5\times9.36=14.0\ \mathrm{V}$（$=4.7\,T_e/e$，符合 $\sim5T_e/e$ 的经验）。$v_B=\sqrt{3\times1.602\times10^{-19}/6.634\times10^{-26}}=2.7\times10^{3}\ \mathrm{m/s}$。}

\item （估算）$n_s=10^{16}\ \mathrm{m^{-3}}$、$T_e=3\ \mathrm{eV}$：求鞘边离子通量 $\Gamma=n_sv_B$ 与德拜长度，并说明为什么这个通量决定了刻蚀速率。
\tkans{$\Gamma=10^{16}\times2.69\times10^{3}=2.7\times10^{19}\ \mathrm{m^{-2}s^{-1}}$（每平方厘米每秒 $2.7\times10^{15}$ 个离子）；$\lambda_D=7430\sqrt{3/10^{16}}=1.3\times10^{-4}\ \mathrm{m}$。每个入射离子带来固定能量与反应概率，故刻蚀/沉积速率正比于 $\Gamma$；提高 $n$ 或 $T_e$ 都能提高通量，但 $T_e$ 同时决定离子能量（鞘电位降）。}

\item （估算）$n=10^{16}\ \mathrm{m^{-3}}$、$T_e=3\ \mathrm{eV}$：等离子体频率 $f_{pe}$ 与趋肤深度 $\delta=c/\omega_{pe}$ 各是多少？为什么 ICP 常用 $13.56\ \mathrm{MHz}$ 而 CCP 对频率不敏感？
\tkans{$\omega_{pe}=5.6\times10^{9}\ \mathrm{rad/s}$，$f_{pe}=9.0\times10^{8}\ \mathrm{Hz}=898\ \mathrm{MHz}$，$\delta=3\times10^{8}/5.6\times10^{9}=5.3\times10^{-2}\ \mathrm{m}$（5.3 cm）。$13.56\ \mathrm{MHz}\ll f_{pe}$，电磁波不能在等离子体中传播，只能以渐逝场（或碰撞/惯性趋肤层）耦合，ICP 靠线圈近场感应耦合，与放电室尺度、频率配合；CCP 靠鞘层电容耦合，离子能量主要由鞘电位降（而非频率）决定，因此对频率不敏感但加热效率受限。}
\end{enumerate}


\section{第 12 章：等离子体诊断}
\tksrc{Chen 第 10 章、Fitzpatrick 第 4 章；探针与波动诊断的严谨处理参照 Hutchinson \textit{Principles of Plasma Diagnostics}。}

\begin{enumerate}
\item （概念）单探针 $I$--$V$ 曲线的三段各对应什么物理过程？如何由曲线求 $T_e$ 与 $n_e$？
\tkans{离子饱和区（$V\ll V_p$）：几乎只收集离子，电流趋于 $I_{is}$；过渡区：电子被部分排斥，电流随 $V$ 指数上升；电子饱和区（$V\gg V_p$）：全部电子被收集，电流趋于 $I_{es}$。$T_e$ 由过渡区 $\ln I$ 对 $V$ 的斜率给出（斜率 $=e/k_BT_e$）；$n_e$ 由 $I_{es}=neA\bar{v}_e/4$（或由 $I_{is}$ 结合玻姆速度）反解。}

\item （概念）汤姆孙散射的 $\alpha$ 判据是什么？什么时候必须用集体散射理论？
\tkans{$\alpha=1/(k\lambda_D)$，$k$ 为散射波矢（$\approx$ 入射与散射波矢之差）。$\alpha\ll1$：散射发生在单个电子的尺度上，各电子独立散射，谱形直接反映电子速度分布（非集体/非相干散射）；$\alpha\gtrsim1$：散射体积内的德拜屏蔽与粒子关联不可忽略，谱中出现离子声波与等离子体波的特征线（集体散射），必须用带动态形状因子的理论解释。$\alpha\sim1$ 附近两种图像都不准，需完整理论。}

\item （推导）由双点法 $I(V_1)/I(V_2)=\exp[e(V_1-V_2)/k_BT_e]$ 给出 $T_e$ 的计算式，并说明取点原则。
\tkans{$T_e=\dfrac{e(V_1-V_2)}{k_B\ln[I(V_1)/I(V_2)]}$。取点须在\textbf{过渡区}且远离浮动电位与饱和区：太靠近 $V_f$ 时离子电流占比大，进入饱和后指数关系失效；同时两点电流比不宜过小（否则 $\ln$ 的相对误差被放大），也不宜跨越探针污染或二次电子发射明显的区域。}

\item （推导）某探针测得 $I(V)=50\ \mathrm{\mu A}$、$I(V+2\ \mathrm{V})=200\ \mathrm{\mu A}$：求 $T_e$；又知电子饱和电流 $I_{es}=1\ \mathrm{mA}$、收集面积 $A=1\ \mathrm{mm^2}$，估算 $n_e$。
\tkans{$T_e=2/\ln4=1.44\ \mathrm{eV}$。$v_{te}=\sqrt{2k_BT_e/m_e}=\sqrt{2\times1.443\times1.602\times10^{-19}/9.109\times10^{-31}}=7.1\times10^{5}\ \mathrm{m/s}$；$n_e=\dfrac{4I_{es}}{eA\bar{v}_e}$（以 $v_{te}$ 近似 $\bar v_e$）$=\dfrac{4\times10^{-3}}{1.602\times10^{-19}\times10^{-6}\times7.1\times10^{5}}=3.5\times10^{16}\ \mathrm{m^{-3}}$。注意用 $v_{te}$ 代替 $\bar v_e$ 会带来约 10\% 的系统偏差（严格用 $\bar v_e=\sqrt{8k_BT_e/\pi m_e}$）。}

\item （估算）微波干涉仪 $f=140\ \mathrm{GHz}$、$\bar n_e=10^{19}\ \mathrm{m^{-3}}$、路径长 $0.5\ \mathrm{m}$：求相移（弧度与条纹数）。
\tkans{截止密度 $n_c=\eps_0m_e\omega^2/e^2=2.43\times10^{20}\ \mathrm{m^{-3}}$；$\lambda=c/f=2.14\times10^{-3}\ \mathrm{m}$；相移 $\Delta\phi=\dfrac{\pi}{\lambda}\dfrac{\bar n_e}{n_c}L=\dfrac{\pi}{2.14\times10^{-3}}\times\dfrac{10^{19}}{2.43\times10^{20}}\times0.5=30.2\ \mathrm{rad}$，即 $30.2/2\pi=4.8$ 条条纹。条纹计数法能分辨约 0.1 条纹，对应的密度分辨约 $2\times10^{17}\ \mathrm{m^{-3}}$。}

\item （估算）磁探针：$N=10$ 匝、$A=10^{-4}\ \mathrm{m^2}$，磁场在 $100\ \mathrm{\mu s}$ 内从 0 升到 $1\ \mathrm{T}$：感应电压是多少？示波器（$\mathrm{mV}$ 灵敏度）能测吗？
\tkans{$\diff B/\diff t=10^{4}\ \mathrm{T/s}$；$V=NA\diff B/\diff t=10\times10^{-4}\times10^{4}=10\ \mathrm{V}$，远高于 $\mathrm{mV}$ 量级，甚至可能过载，实际探针常减小 $N$、$A$ 或加衰减。反向使用时，把 $V$ 数值积分即可还原 $B(t)$，这正是磁探针（及 Mirnov 线圈阵列）测磁涨落与 MHD 模的基本手段。}
\end{enumerate}


\section{来源、分级与版权说明}
\label{sec:app-l-source}

\begingroup\small
本附录的选题原则是\textbf{借鉴题型与分级，不复制题目}。下表逐章说明每题型的参照来源；表后是版权与独立性声明。

{\small
\begin{longtable}{@{}p{1.4cm}p{5.3cm}p{5.4cm}p{3.4cm}@{}}
\caption{分章题库的题型来源与难度参照（与附录 J 的逐章对照表 \ref{tab:app-j-map} 配套使用）}\label{tab:app-l-source}\\
\toprule
本书章 & 借鉴的教材与章节 & 借鉴的题型（不是题目） & 本题库构成 \\
\midrule
\endfirsthead
\toprule
本书章 & 借鉴的教材与章节 & 借鉴的题型（不是题目） & 本题库构成 \\
\midrule
\endhead
\bottomrule
\endlastfoot
1 & Chen 第 1 章及附录 A；Fitzpatrick 各章数学工具 & 分布函数积分、单位与量纲换算、屏蔽长度的量级题 & 概念 2 / 推导 2 / 估算 2 \\
2 & Chen 第 1 章；Fitzpatrick 第 1 章 & 判据的边界讨论、德拜球内粒子数估算 & 概念 2 / 推导 2 / 估算 2 \\
3 & Chen 第 2 章；Fitzpatrick 第 2 章 & 回旋频率与半径的数值题、损失锥几何 & 概念 2 / 推导 2 / 估算 2 \\
4 & Chen 第 3 章；Fitzpatrick 第 4 章 & 流体方程的守恒性讨论、抗磁漂移推导 & 概念 2 / 推导 2 / 估算 2 \\
5 & Chen 第 3、6 章；Fitzpatrick 第 8 章 & Alfvén 速度、磁扩散时间、$\beta$ 的标度题 & 概念 2 / 推导 2 / 估算 2 \\
6 & Chen 第 4 章；Fitzpatrick 第 5、6、7 章 & 截止密度、色散关系的极限讨论、波的传播时间 & 概念 2 / 推导 2 / 估算 2 \\
7 & Chen 第 7 章；Fitzpatrick 第 3、7 章 & 线性化推导、弱阻尼公式的数值代入与敏感度 & 概念 2 / 推导 2 / 估算 2 \\
8 & Chen 第 5 章；Fitzpatrick 第 3、4 章 & 库仑对数与碰撞时间、平均自由程与输运标度 & 概念 2 / 推导 2 / 估算 2 \\
9 & Chen 第 6、8 章；Fitzpatrick 第 7、9 章 & 增长率推导、瑞利--泰勒与交换模的类比题 & 概念 2 / 推导 2 / 估算 2 \\
10 & Chen 第 9、10 章；Goldston \& Rutherford；Wesson & 零维功率平衡、点火门槛、堆芯功率与体积估算 & 概念 2 / 推导 2 / 估算 2 \\
11 & Chen 第 10 章；Fitzpatrick 第 4 章末；Lieberman \& Lichtenberg & 鞘层与玻姆判据、浮动电位、通量与耦合方式 & 概念 2 / 推导 2 / 估算 2 \\
12 & Chen 第 10 章；Hutchinson & 探针特性曲线的读数、干涉仪相移、磁探针感应电压 & 概念 2 / 推导 2 / 估算 2 \\
\end{longtable}
}

\textbf{版权与独立性声明：}
\begin{itemize}
\item 上表所列「借鉴」仅指\textbf{题型分类与难度定位}（例如「求德拜长度的量级估算」「由玻姆判据讨论鞘层」这类公认的练习形式），属于学科共同的教学传统；
\item 本附录全部题目的\textbf{文字表述、参数选取与数值答案均由本书自拟}，并用仓库里的校核脚本逐条复算（两位有效数字）；
\item 本附录\textbf{不复制}任何教材的题目原文、解答、图表或数据表；引用教材仅给出书名与章节，便于读者回原书查找（完整书目见附录 E 与附录 J）；
\item 若读者发现某题与某教材的题目高度雷同，属独立命制的巧合，欢迎指出，我们会在下一版改写或补注出处。
\end{itemize}

\begin{tip}{做完 72 题之后}{tip-appl-next}
如果你把 12 章的 72 道题都做完了，可以回到附录 I 的第 5 节，把这 12 周的计划换成「每周一章：6 题自测 + 1 道附录 K 的范例精读」，题库与范例一一对应，正好覆盖全书主线。
\end{tip}
\endgroup
